(1) Fraunhofer振幅正比$\int_b^a r\,dr\int_0^{2\pi}e^{-ikr\sin\theta\cos\phi}\,d\phi=2\pi\int_b^a rJ_0(ur/a)\,dr$。用恒等式$\frac{d}{dr}[rJ_1(qr)]=qrJ_0(qr)$积分，得到$F(u)\propto (2\pi a^2/u)[J_1(u)-\epsilon J_1(\epsilon u)]$，这是两个同轴圆孔的振幅相减，而不是强度相减。
(2) 轴向极限$J_1(x)=x/2+O(x^3)$使$F(0)\propto\pi a^2(1-\epsilon^2)$，正好等于透光环面积。在单位轴上强度归一化下，$A(u)=2[J_1(u)-\epsilon J_1(\epsilon u)]/[u(1-\epsilon^2)]$，相对强度$I(u)/I(0)=|A(u)|^2$。对固定照明幅度、同一传播距离，轴上绝对强度随$(1-\epsilon^2)^2$下降。
(3) 非轴向第一零点由$J_1(u)=\epsilon J_1(\epsilon u)$给出，并应取最小正根（不能把$u=0$误作暗环）。若$\epsilon=0$，方程变为$J_1(u)=0$，第一根为$u_1\simeq3.831706$；因此圆孔第一暗环角满足$\sin\theta_1\simeq3.831706/(ka)=1.22\lambda/(2a)$，与标准Airy衍射一致。
(4) 当$\epsilon\to1^-$，对$g(\epsilon)=\epsilon J_1(\epsilon u)$作Taylor展开，$J_1(u)-g(\epsilon)\sim(1-\epsilon)[J_1(u)+uJ_1'(u)]=(1-\epsilon)uJ_0(u)$。分母$u(1-\epsilon^2)\sim2u(1-\epsilon)$，于是$A(u)\to J_0(u)$。第一暗纹趋$J_0$第一正根$u_1\simeq2.404826$，比普通圆孔更靠近中心，但副瓣能量分布会增强；仅凭第一零点变小不能断言总体成像对比度更优。
